\part{How to Use the Memory Workbook}

\chapter{Recall by Study Event}

This workbook is the cumulative memory companion to the twenty-two-packet
common trunk and the eight-packet advanced poetry branch. Its intervals are
\emph{study events}: a return happens after later packets have been studied,
however much time that takes. Repeat any sheet as often as useful. Write from
memory before checking the model, correct visibly, state the reason for each
substantive correction, and carry uncertain material forward.

The controlled registry contains \arabic{MemoryLexemeCount} lexemes and
\arabic{MemoryGrammarFactCount} grammar facts allocated across
\arabic{MemorySetCount} packet sets.  Every packet provides the same eight
printable sheet families:

\begin{xltabular}{\linewidth}{@{}P{1.2in}P{2.0in}Y@{}}
  \toprule
  \textbf{Suffix} & \textbf{Sheet} & \textbf{Work}\\
  \midrule
  \texttt{-S} & New-lexeme study & Complete citation form, project gloss,
    morphology, usage or government, and pedagogical strand\\
  \texttt{-L2E} & Latin to sense & New words plus returns after one, two, and
    five later packets and at selected stage boundaries\\
  \texttt{-E2L} & Sense to Latin & Complete citation-form reconstruction in
    the reverse direction\\
  \texttt{-MORPH} & Morphology & Declension, gender, adjective pattern,
    principal parts, conjugation, or indeclinable class\\
  \texttt{-GOV} & Usage and government & Case government, complement pattern,
    construction, or lexical warning, followed by an original sentence\\
  \texttt{-GRAM} & Grammar facts & Rule, diagnostic, distinction, or method
    stated in complete grammatical language\\
  \texttt{-RECON} & Reconstruction & Ordered Latin, complete parsing, and an
    account of every supplied form\\
  \texttt{-COMP} & Composition & Proposition outline, checked draft, and
    revision on the way to advanced prose or verse composition\\
  \bottomrule
\end{xltabular}

Thus \texttt{M01-L2E} always names the Latin-to-sense sheet for packet M01,
and \texttt{P8-COMP} always names the composition sheet at the end of the
poetry branch.  These stable IDs make a useful personal return list: record the
exact sheet and lexeme or grammar-fact IDs that remain uncertain.

\section{The five recurrence labels}

\begin{xltabular}{\linewidth}{@{}P{0.65in}P{1.45in}Y@{}}
  \toprule
  \textbf{Label} & \textbf{Event} & \textbf{Use}\\
  \midrule
  N & New & First deliberate study and immediate retrieval in the packet where
    the item enters the controlled memory set\\
  R1 & One-packet return & Retrieval after one later packet has intervened\\
  R2 & Two-packet return & Retrieval after two later packets have intervened\\
  R5 & Five-packet return & Wider retrieval after five later packets have
    intervened\\
  S & Stage return & Deliberately mixed cumulative return near a structural
    boundary or the close of a branch\\
  \bottomrule
\end{xltabular}

The schedule is a floor, not a ceiling.  A difficult item can appear on a
personal return list at every sitting until it is dependable.  An easy item
still receives the printed widening returns, because fluent reading depends on
rapid availability, not one successful recognition.

\section{Source and exercise boundary}

The English senses in this workbook are concise project-created study glosses.
They do not replace a full dictionary entry.  A label such as ``general and
classical prose'' or ``biblical and liturgical reading'' assigns pedagogical
work; it does not claim that the registry entry has been collated at a named
author, work, edition, and locus.  Exact received wording belongs in the
course's passage inventory and source records after verification.  Every Latin
reconstruction and composition prompt or model in this workbook is
project-created instructional material, not an unattributed quotation.

\section{Before continuing}

It is reasonable to continue when the packet's citation forms, principal
senses, essential morphological facts, and constructions can usually be
recalled from both Latin and English cues; when its grammar prompts can be
explained rather than merely recognized; and when the same material remains
available in fresh connected Latin and short composition.  Dependable recall
matters more than perfect recall.  Record what is uncertain and keep it in the
return cycle while later reading supplies further encounters.

\part{Common-Trunk Practice Sheets}

\RenderMemoryWorkbookLearnerPacket{M01}
\RenderMemoryWorkbookLearnerPacket{M02}
\RenderMemoryWorkbookLearnerPacket{M03}
\RenderMemoryWorkbookLearnerPacket{M04}
\RenderMemoryWorkbookLearnerPacket{M05}
\RenderMemoryWorkbookLearnerPacket{M06}
\RenderMemoryWorkbookLearnerPacket{M07}
\RenderMemoryWorkbookLearnerPacket{M08}
\RenderMemoryWorkbookLearnerPacket{M09}
\RenderMemoryWorkbookLearnerPacket{M10}
\RenderMemoryWorkbookLearnerPacket{M11}
\RenderMemoryWorkbookLearnerPacket{M12}
\RenderMemoryWorkbookLearnerPacket{M13}
\RenderMemoryWorkbookLearnerPacket{M14}
\RenderMemoryWorkbookLearnerPacket{M15}
\RenderMemoryWorkbookLearnerPacket{M16}
\RenderMemoryWorkbookLearnerPacket{M17}
\RenderMemoryWorkbookLearnerPacket{M18}
\RenderMemoryWorkbookLearnerPacket{M19}
\RenderMemoryWorkbookLearnerPacket{M20}
\RenderMemoryWorkbookLearnerPacket{M21}
\RenderMemoryWorkbookLearnerPacket{M22}

\part{Advanced Poetry-Branch Practice Sheets}

Poetry is prepared within the common trunk through close attention to syntax,
word order, compression, rhetoric, textual evidence, and composition.  The
following branch adds explicit prosodic, metrical, poetic-syntactic, and verse-
composition memory work; it does not displace continued Missal and prose
reading.

\RenderMemoryWorkbookLearnerPacket{P1}
\RenderMemoryWorkbookLearnerPacket{P2}
\RenderMemoryWorkbookLearnerPacket{P3}
\RenderMemoryWorkbookLearnerPacket{P4}
\RenderMemoryWorkbookLearnerPacket{P5}
\RenderMemoryWorkbookLearnerPacket{P6}
\RenderMemoryWorkbookLearnerPacket{P7}
\RenderMemoryWorkbookLearnerPacket{P8}

\part{Common-Trunk Models and Checks}

Use these pages after a fresh attempt.  A model supplies one concise answer;
where context permits more than one defensible sense, construction, or
composition, compare the grammatical evidence rather than treating different
wording as an error by itself.

\RenderMemoryWorkbookSolutionPacket{M01}
\RenderMemoryWorkbookSolutionPacket{M02}
\RenderMemoryWorkbookSolutionPacket{M03}
\RenderMemoryWorkbookSolutionPacket{M04}
\RenderMemoryWorkbookSolutionPacket{M05}
\RenderMemoryWorkbookSolutionPacket{M06}
\RenderMemoryWorkbookSolutionPacket{M07}
\RenderMemoryWorkbookSolutionPacket{M08}
\RenderMemoryWorkbookSolutionPacket{M09}
\RenderMemoryWorkbookSolutionPacket{M10}
\RenderMemoryWorkbookSolutionPacket{M11}
\RenderMemoryWorkbookSolutionPacket{M12}
\RenderMemoryWorkbookSolutionPacket{M13}
\RenderMemoryWorkbookSolutionPacket{M14}
\RenderMemoryWorkbookSolutionPacket{M15}
\RenderMemoryWorkbookSolutionPacket{M16}
\RenderMemoryWorkbookSolutionPacket{M17}
\RenderMemoryWorkbookSolutionPacket{M18}
\RenderMemoryWorkbookSolutionPacket{M19}
\RenderMemoryWorkbookSolutionPacket{M20}
\RenderMemoryWorkbookSolutionPacket{M21}
\RenderMemoryWorkbookSolutionPacket{M22}

\part{Advanced Poetry-Branch Models and Checks}

\RenderMemoryWorkbookSolutionPacket{P1}
\RenderMemoryWorkbookSolutionPacket{P2}
\RenderMemoryWorkbookSolutionPacket{P3}
\RenderMemoryWorkbookSolutionPacket{P4}
\RenderMemoryWorkbookSolutionPacket{P5}
\RenderMemoryWorkbookSolutionPacket{P6}
\RenderMemoryWorkbookSolutionPacket{P7}
\RenderMemoryWorkbookSolutionPacket{P8}
